% cap05/5.4_roi.tex — Impacto transaccional y ROI arquitectonico.
\section{Impacto Transaccional y Retorno de Inversi\'on (ROI) Arquitect\'onico}

M\'as all\'a de los vectores puramente computacionales y psicom\'etricos, el proyecto debe tributar a las ciencias econ\'omicas institucionales.

\subsection{Reducci\'on de Tiempos de Consolidaci\'on de Datos}
En un sondeo transaccional operativo, se midi\'o el tiempo (\textit{Time to task completion}) que un administrador financiero t\'ipico de ONG invert\'ia en reconciliar 50 donaciones y asignarlas a 30 voluntarios. Mediante hojas de c\'alculo (m\'etodo pre-existente en su organizaci\'on l\'inea base), el proceso de consolidaci\'on devoraba, en promedio, $145 \text{ minutos}$.
Al desplegar las funciones reactivas automatizadas del panel \textit{Democra}, el mismo experimento baj\'o asint\'oticamente a $18 \text{ minutos}$. Esto implica una optimizaci\'on de flujo de trabajo humano y reducci\'on de sobrecarga administrativa de un 87.5\%.

Estas horas-hombre rescatadas de la burocracia son horas que el personal de la ONG podr\'a re-injetar directamente en sus l\'ineas operativas de campo, demostrando as\'i que la ingenier\'ia de sistemas y las arquitecturas Multi-Tenant de alto nivel resultan cr\'iticamente salvadoras en escenarios humanitarios.
